\section{The \texttt{Virtual\_input} McStas Component}
Source-like component that generates neutron events from an ascii
'virtual source' filename.

\subsection*{Identification}
\begin{itemize}
  \item \textbf{Site:} 
  \item \textbf{Author:} \textless{}a href="mailto:farhi@ill.fr"\textgreater{}E. Farhi\textless{}/a\textgreater{}
  \item \textbf{Origin:} \textless{}a href="http://www.ill.fr"\textgreater{}ILL\textless{}/a\textgreater{}
  \item \textbf{Date:} Sep 28th, 2001
\end{itemize}

\subsection*{Description}
\begin{lstlisting}
This component reads neutron events stored from a file, and sends them into
the instrument. It thus replaces a Source component, using a previously
computed neutron set. The 'source' file type is an ascii text file with the
format listed below. The number of neutron events for the
simulation is set to the length of the 'source' file times the
repetition parameter 'repeat_count' (1 by default).
It is particularly useful to generate a virtual source at a point that few
neutron reach. A long simulation will then only be performed once, to create
the 'source' filename. Further simulations are much faster if they start from
this low flux position with the 'source' filename.

The input file format is:
text column formatted with lines containing 11 values in the order:
p x y z vx vy vz t sx sy sz stored into about 83 bytes/n.

%BUGS
We recommend NOT to use parallel execution (MPI) with this component. If you
need to, set parameter 'smooth=1'.

EXAMPLE:
To create a 'source' file collecting all neutron states, use:
COMPONENT MySourceCreator = Virtual_output(filename = "MySource.list")
at the position where will be the Virtual_input.
Then inactivate the part of the simulation description before (and including)
the component MySourceCreator. Put the new instrument source:
COMPONENT Source = Virtual_input(filename="MySource.list")
at the same position as 'MySourceCreator'.
A Vitess filename may be obtained from the 'Vitess_output' component or from a
Vitess simulation (104 bytes per neutron) and read with Vitess_input.
\end{lstlisting}

\subsection*{Input parameters}
Parameters in \textbf{boldface} are required; the others are optional.

\begin{longtable}{p{0.22\textwidth}p{0.12\textwidth}p{0.46\textwidth}p{0.14\textwidth}}
\toprule
\textbf{Name} & \textbf{Unit} & \textbf{Description} & \textbf{Default} \\
\midrule
\endhead
filename & str & Name of the neutron input file. Empty string "" inactivates component & 0 \\
verbose & 0/1 & Display additional information about source, recommended & 0 \\
repeat\_count & 1 & Number of times the source must be generated/repeated & 1 \\
smooth & 0/1 & Smooth sparsed event files for filename repetitions. Use this option with MPI. This will apply gaussian distributions around initial events from the file & 1 \\
\bottomrule
\end{longtable}

\subsection*{Links}
\begin{itemize}
  \item \href{run:/home/willend/willend-McCode/mcstas-comps/obsolete/Virtual_input.comp}{Source code} for \texttt{Virtual\_input.comp}.
\end{itemize}
\IfFileExists{Virtual_input_static.tex}{\input{Virtual_input_static.tex}}{}